% AUTHOR: Amlal El Mahrouss
% PURPOSE: A Short Reference for Algebraic Properties.

\documentclass[10pt]{article}

%% ---- Core Packages ----------------------------------------
\usepackage[T1]{fontenc}
\usepackage[utf8]{inputenc}
\usepackage{lmodern}
\usepackage{microtype}

%% ---- Page Geometry ----------------------------------------
\usepackage[
top=1in, bottom=1in,
left=0.75in, right=0.75in,
columnsep=0.25in
]{geometry}

%% ---- Mathematics ------------------------------------------
\usepackage{amsmath, amssymb, amsthm, amsfonts}
\usepackage{mathtools}
\usepackage{bm}
\usepackage{textcomp}
\usepackage{esint}

%% ---- Figures & Tables -------------------------------------
\usepackage{graphicx}
\usepackage{booktabs}
\usepackage{multirow}
\usepackage{array}
\usepackage{caption}
\usepackage{subcaption}
\usepackage{float}

%% ---- Code Listings ----------------------------------------
\usepackage{listings}
\usepackage{xcolor}
\usepackage{algorithm}
\usepackage{algpseudocode}

\lstdefinestyle{codestyle}{
	backgroundcolor=\color{gray!8},
	basicstyle=\ttfamily\footnotesize,
	breakatwhitespace=false,
	breaklines=true,
	captionpos=b,
	commentstyle=\color{green!50!black},
	keywordstyle=\color{blue}\bfseries,
	stringstyle=\color{orange!70!black},
	numberstyle=\tiny\color{gray},
	numbers=left,
	numbersep=5pt,
	showstringspaces=false,
	frame=single,
	rulecolor=\color{gray!40},
	tabsize=2
}
\lstset{style=codestyle}

%% ---- Hyperlinks -------------------------------------------
\usepackage[
colorlinks=true,
linkcolor=blue!70!black,
citecolor=green!50!black,
urlcolor=blue!60!black
]{hyperref}
\usepackage{url}

%% ---- Bibliography -----------------------------------------
\usepackage[numbers, sort&compress]{natbib}

%% ---- Miscellaneous ----------------------------------------
\usepackage{enumitem}
\usepackage{siunitx}
\usepackage{cleveref}

%% ---- Theorem Environments ---------------------------------
\theoremstyle{plain}
\newtheorem{theorem}{Theorem}[section]
\newtheorem{lemma}[theorem]{Lemma}
\newtheorem{corollary}[theorem]{Corollary}
\newtheorem{proposition}[theorem]{Proposition}

\theoremstyle{definition}
\newtheorem{definition}[theorem]{Definition}
\newtheorem{example}[theorem]{Example}

\theoremstyle{remark}
\newtheorem{remark}[theorem]{Remark}

%% ---- Custom Commands --------------------------------------
\newcommand{\R}{\mathbb{R}}
\newcommand{\N}{\mathbb{N}}
\newcommand{\Z}{\mathbb{Z}}
\newcommand{\E}{\mathbb{E}}
\newcommand{\Prob}{\mathbb{P}}
\newcommand{\norm}[1]{\left\lVert #1 \right\rVert}
\newcommand{\abs}[1]{\left\lvert #1 \right\rvert}
\newcommand{\inner}[2]{\left\langle #1,\, #2 \right\rangle}
\newcommand{\ie}{\textit{i.e.}\xspace}
\newcommand{\eg}{\textit{e.g.}\xspace}
\newcommand{\etal}{\textit{et al.}\xspace}

%% ---- Caption Formatting -----------------------------------
\captionsetup{
	font=small,
	labelfont=bf,
	format=hang,
	justification=justified
}

\renewcommand{\qedsymbol}{\ensuremath{\blacksquare}}

\usepackage{lettrine}

%% ---- Header / Footer --------------------------------------
\usepackage{fancyhdr}
\pagestyle{fancy}
\fancyhf{}
\fancyhead[R]{\small\thepage}
\renewcommand{\headrulewidth}{0.4pt}

%% ============================================================
%%  FRONT MATTER
%% ============================================================

\title{%
	\vspace{-1.5em}%
	\large\textbf{A Short Reference for Algebraic Properties.}%
}
\author{By A. E. Mahrouss}
\date{\today}

%% ============================================================
\begin{document}
%% ============================================================

\maketitle
\tableofcontents
\newpage

%% ============================================================
\section{Synopsis:}
%% ============================================================

\lettrine[lines=1, lhang=0.1, loversize=0.2]{L}ET the following theorems, proof, and definitions for this note.

%% ============================================================
\section{Observations:}
%% ============================================================

\lettrine[lines=1, lhang=0.1, loversize=0.2]{L}ET the following observations:

\begin{theorem}
	There exists for $p, \operatorname{L}, f(a, qk), a, q, k \in \mathbb{R}.$:
	\begin{equation}
		\lim_{p \to \operatorname{L}} (\frac{f(a, qk)}{k} + (\pi^{1/12} - \frac{\zeta(3)}{\pi}))^p, \quad f(a, qk), p \neq 0.
	\end{equation}
	Solutions of $\lim_{p \to \operatorname{L}} (\frac{f(a, qk)}{k} + (\pi^{1/12} - \frac{\zeta(3)}{\pi}))^p,$ in  $\operatorname{L} \in \{2k : k \in \mathbb{Z}\}.$
\end{theorem}
\begin{remark}
	There exists an infinite set of solutions in $\operatorname{L} \in \{2k : k \in \mathbb{Z}\}$ in parity.
\end{remark}
\begin{remark}
There exists an infinite set of solutions of transcendental property.
\end{remark}

\begin{remark}
	Per recent exploration of properties, It may be worth sharing this specific bounded integral:
	\begin{align*}
		\int_{\frac{1}{2}}^{2} (\pi^{3/12} + \frac{\zeta(2)}{\pi}) &= \frac{1}{4} (6 \cdot \pi^{\frac{1}{4}} + \pi),\\
		\int_{\frac{1}{2}}^{2} (\pi^{3/12} + \frac{\zeta(2)}{\pi}) &= \frac{1}{4} (6 \cdot \sqrt[4]{\pi} + \pi).
	\end{align*}
	May admit a sequence of solutions in $\mathbb{R}$, of algebraic and transcendence property. Posing the limit of $(\pi^{3/12} + \frac{\zeta(2)}{\pi})$ in $+\infty$, we get:
	\begin{align*}
	\lim_{x \to 0} (\pi^{3/12} + \frac{\zeta(2)}{\pi}) &= \frac{1}{4} (6 \cdot \pi^{\frac{1}{4}} + \pi).
	\end{align*}
	Which seems to be generalizable under the following term for $n \neq 0, \quad n \in \mathbb{N}.$ And a non-zero $p \in \mathbb{R}.$:
	\begin{align*}
		\int_{\frac{1}{p}}^{p} (\pi^{3/12} + \frac{\zeta(n)}{\pi}).
	\end{align*}
\end{remark}

\begin{corollary}
	Per Remark 2.1 for $p \neq 0, \frac{1}{p} < \zeta(n)$, per Remark 2.1, $\int_{\frac{1}{p}}^{p} (\pi^{3/12} + \frac{\zeta(n)}{\pi}) \neq \frac{1}{p}$, thus $\frac{1}{p} < \int_{\frac{1}{p}}^{p} (\pi^{3/12} + \frac{\zeta(n)}{\pi}).$
\end{corollary}

\begin{corollary}
	Per Corollary 3.1, $\int_{\frac{1}{p}}^{p} (\pi^{3/12} + \frac{\zeta(n)}{\pi})$, thus $\int_{\frac{1}{p}}^{p} (\pi^{3/12} + \frac{\zeta(n)}{\pi}) \geq \frac{1}{p}$. And $\int_{\frac{1}{p}}^{p} (\pi^{3/12} + \frac{\zeta(n)}{\pi}) \leq p$, thus $\frac{1}{p} \leq \int_{\frac{1}{p}}^{p} (\pi^{3/12} + \frac{\zeta(n)}{\pi}) \leq p$.
\end{corollary}

\begin{corollary}
	There exists an indetity of the following for $\forall t, k, x \in \mathbb{R}$.:
	\begin{align*}
		\mathcal{M}(x, k) = \int_{n}^{p} \mathcal{M}_{x}(\mathcal{X}_k) \cdot t \quad dt,\\
		\mathcal{M}(x, k)^{2} = \int_{n}^{p} \mathcal{M}_{x}(\mathcal{X}_k) \cdot t \quad dt^{2}.
	\end{align*}
	And thus we may write the following:
	\begin{align*}
		\mathcal{M}(x, k)^{2} + \mathcal{M}(x, k).
	\end{align*}
	Which then could be rewritten as:
	\begin{align*}
		F(x, k) &= \mathcal{M}(x, k)^{2} + \mathcal{M}(x, k).
	\end{align*}
	Adding the tensors defined as: $\mathcal{T}_3, \mathcal{T}_2 \in \mathbb{R}.$:
	\begin{align*}
		F(x, k) + \mathcal{T}_2 &= \mathcal{M}(x, k)^{2} + \mathcal{M}(x, k) + \mathcal{T}_3.
	\end{align*}
	A more complete definition with $\mathcal{T}_2$ and $\mathcal{T}_3$ would then be:
	\begin{align*}
		\mathcal{T}_2 &= \mathcal{M}(x, k)^{2} + \mathcal{M}(x, k)^{1} + \mathcal{T}_3 - F(x, k), \quad \forall x, k \in \mathbb{R}.
	\end{align*}
	A more rigorous form of the equation above would be:
	\begin{align*}
			\pm \mathcal{T}_2 &= \mathcal{M}(x, k)^{2} + \mathcal{M}(x, k)^{1} + \mathcal{T}_3 \pm F(x, k), \quad \forall x, k \in \mathbb{R}.
	\end{align*}
\end{corollary}

%% ============================================================
\section{Going Forward:}
%% ============================================================

\lettrine[lines=1, lhang=0.1, loversize=0.2]{N}OW, we may start working on further notes regarding riemannian and differential geometry, alongside several advanced algebra concepts.

\nocite{*}

\bibliographystyle{acm}
\bibliography{chambadal1981dictionnaire, pinter2010book, Riemann, fatou1906, weil1967basic, A_Course_of_Pure_Mathematics}

%% ============================================================
\end{document}
%% ============================================================
